% Generated by ManuScript from .ms/approved/draft-approved-v1.md (sha256 e9e1bd087f4c4544). Edit the Markdown, not this file, unless you accept re-approval.
\documentclass[conference]{IEEEtran}
\IEEEoverridecommandlockouts
\usepackage{cite}
\usepackage{amsmath,amssymb,amsfonts}
\usepackage{algorithmic}
\usepackage{graphicx}
\usepackage{textcomp}
\usepackage{xcolor}
\usepackage{fancyvrb}
\usepackage{url}
\def\BibTeX{{\rm B\kern-.05em{\sc i\kern-.025em b}\kern-.08em
    T\kern-.1667em\lower.7ex\hbox{E}\kern-.125emX}}
\newcommand{\pandocbounded}[1]{#1}
\providecommand{\tightlist}{\setlength{\itemsep}{0pt}\setlength{\parskip}{0pt}}
\setkeys{Gin}{width=\linewidth,height=0.4\textheight,keepaspectratio}
\providecommand{\passthrough}[1]{#1}
\usepackage[hidelinks]{hyperref}
\begin{document}

\title{A Fixture Report}

\author{\IEEEauthorblockN{Test Author}}

\maketitle

\begin{abstract}
This is a test abstract.
\end{abstract}


\section{Introduction}\label{introduction}

Prior work \cite{smith2020} motivates this study. See
Fig.~\ref{fig-arch} and Table~\ref{tbl-res}.

\section{Methodology}\label{methodology}

We follow the approach of \cite{lee2018} and describe it here.

\begin{figure}
\centering
\pandocbounded{\includegraphics[keepaspectratio,alt={System
architecture}]{figures/arch.png}}
\caption{System architecture}\label{fig-arch}
\end{figure}

\section{Results}\label{results}

\begin{table}[htbp]
\caption{Results summary}\label{tbl-res}
\centering\footnotesize
\begin{tabular}{l l}
\hline
\textbf{Metric} & \textbf{Value} \\
\hline
A & 1 \\
B & 2 \\
\hline
\end{tabular}
\end{table}

\section{Conclusion}\label{conclusion}

The work is complete \cite{lee2018}.

\bibliographystyle{IEEEtran}
\bibliography{references}

\end{document}
